\chapter*{Conclusioni}
\addcontentsline{toc}{chapter}{Conclusioni}

\section*{Considerazioni finali}

Il presente lavoro non ha riscontrato, sui tre dataset reali analizzati e nelle condizioni sperimentali adottate, alcun vantaggio predittivo dei modelli di Quantum Machine Learning rispetto ai corrispondenti modelli classici di riferimento; anzi, il divario di prestazioni a favore di questi ultimi si è mostrato sistematico e crescente con la complessità del problema. Si tratta di un risultato negativo ma non casuale, poiché riconducibile a meccanismi specifici individuati nel Capitolo~\ref{cap:implementazione_sperimentale_risultati}: la difficoltà di ottimizzazione del VQC in presenza di barren plateaus, la concentrazione del kernel quantistico all'aumentare del numero di qubit e la sensibilità al rumore proporzionale alla profondità del circuito. Il rigore del protocollo sperimentale adottato, con baseline classiche ottimizzate tramite ricerca a griglia e un confronto a parità di dati in ingresso, rende questo esito una misura affidabile dei limiti attuali del QML su problemi di classificazione di questa scala, e non un artefatto di un confronto impari.

\section*{Prospettive future}

I risultati e i limiti discussi nel Capitolo~\ref{cap:implementazione_sperimentale_risultati} suggeriscono alcune direzioni prioritarie per lavori futuri. La più immediata è il completamento della campagna sperimentale su hardware quantistico reale, pianificando per tempo l'accesso a un piano di calcolo adeguato o distribuendo l'esecuzione su una finestra temporale compatibile con la quota gratuita disponibile, in modo da validare empiricamente, su dispositivo fisico, il degrado di prestazioni finora osservato solo in simulazione rumorosa.

Sarebbe inoltre utile estendere l'analisi di sensitività condotta sul dataset Digits (Sezione~\ref{sec:risultati_digits}) ad architetture di feature map e ansatz alternative a quelle qui adottate (ZZFeatureMap e RealAmplitudes), per verificare se il fenomeno di concentrazione del kernel osservato sia una caratteristica intrinseca della codifica scelta o un effetto più generale.
Un altro approccio consiste nel condurre una ricerca, utilizzando gli stessi dataset o altri, sulla struttura di correlazione tra le feature più difficili da catturare con un kernel classico: la letteratura scientifica riconosce infatti la presenza di tale struttura come condizione necessaria, ma non sufficiente, per l'esistenza di un vantaggio quantistico.
Infine l'evoluzione dell'hardware quantistico verso un numero maggiore di qubit e una migliore fedeltà delle porte logiche potrebbe in futuro consentire di ripetere questo confronto su dataset di dimensioni maggiori, con riduzioni dimensionali diverse.
